\documentclass[11pt,a4paper]{article}

\usepackage[utf8]{inputenc}
\usepackage[T1]{fontenc}
\usepackage{lmodern}
\usepackage{microtype}
\usepackage[english]{babel}
\usepackage[a4paper,margin=2.5cm,headheight=14pt]{geometry}
\usepackage{xcolor}
\usepackage{listings}
\usepackage{enumitem}
\usepackage{booktabs}
\usepackage{fancyhdr}
\usepackage{hyperref}
\usepackage{titlesec}

\hypersetup{
    colorlinks=true,
    linkcolor=black,
    urlcolor=blue!60!black,
    citecolor=black,
}

\definecolor{codebg}{RGB}{248,248,242}
\definecolor{codecomment}{RGB}{106,135,89}
\definecolor{codekeyword}{RGB}{0,102,153}
\definecolor{codestring}{RGB}{170,55,55}
\definecolor{errred}{RGB}{200,40,40}

\lstdefinestyle{recipix}{
    backgroundcolor=\color{codebg},
    basicstyle=\ttfamily\small,
    keywordstyle=\color{codekeyword}\bfseries,
    commentstyle=\color{codecomment}\itshape,
    stringstyle=\color{codestring},
    breaklines=true,
    showstringspaces=false,
    frame=single,
    framerule=0pt,
    framesep=4pt,
    aboveskip=0.6em,
    belowskip=0.6em,
    morekeywords={recipe,function,ingredient,step,let,if,else,
        repeat,times,foreach,in,at,for,evaluate,scale,substitute,
        serves,with,by,ratio,return,quantity\_of,true,false,
        int,float,bool,Mass,Volume,Count,Temperature,Duration,Pinch,
        combine,mix,pour,melt,whisk,blend,bake,flip,add,sprinkle,drizzle,knead},
    morecomment=[l]{//},
    morecomment=[s]{/*}{*/},
    morestring=[b]",
}

\lstdefinestyle{output}{
    backgroundcolor=\color{codebg},
    basicstyle=\ttfamily\small,
    breaklines=true,
    showstringspaces=false,
    frame=single,
    framerule=0pt,
    framesep=4pt,
    aboveskip=0.6em,
    belowskip=0.6em,
}

\lstdefinestyle{errmsg}{
    backgroundcolor=\color{codebg},
    basicstyle=\ttfamily\small\color{errred},
    breaklines=true,
    showstringspaces=false,
    frame=single,
    framerule=0pt,
    framesep=4pt,
    aboveskip=0.4em,
    belowskip=0.4em,
}

\lstset{style=recipix}

\titleformat{\section}{\Large\bfseries}{\thesection.}{0.6em}{}
\titleformat{\subsection}{\large\bfseries}{\thesubsection}{0.6em}{}

\pagestyle{fancy}
\fancyhf{}
\fancyhead[L]{CSE 341 -- Recipix v4.1}
\fancyhead[R]{D3 (Part 2)}
\fancyfoot[C]{\thepage}

\title{\textbf{D3 -- Test Report (Part 2)} \\
       \large Recipix v4.1 -- sample programs, outputs, and type-error trace}
\author{Berk Hakan \"Oge (210104004132) \and \.Ismail Kabak (1901042652)}
\date{22 May 2026}

\begin{document}
\maketitle
\thispagestyle{fancy}


\section{Overview}

This report documents the Part~2 test artefacts for Recipix v4.1. Per
handout \S D3~[P2], we provide:
\begin{itemize}[leftmargin=*,itemsep=2pt]
  \item \textbf{Three non-trivial example programs} that together exercise
        every control structure, the structured type (recipes and
        ingredients), user-defined functions and calls, and the three
        domain-specific operators (\texttt{evaluate}, \texttt{scale},
        \texttt{substitute}).
  \item \textbf{Actual rendered outputs} from the interpreter for the two
        well-typed samples, with 3--5 sentence discussion paragraphs each.
  \item \textbf{One type-error program} that exercises the type checker
        and is rejected before execution, with the verbatim error message
        the checker produces.
  \item \textbf{Five parse-error fixtures} carried over from Part~1, with
        the parser error messages they produce.
\end{itemize}

All programs live in \texttt{src/tests/fixtures/} and are reproducible
from the repo. To run a program end-to-end:

\begin{lstlisting}[style=output]
# Parse only (always works):
python main.py src/tests/fixtures/valid/sample1.rcx

# Type-check only (uses the Part 2 type checker):
python main.py --typecheck src/tests/fixtures/valid/sample3.rcx

# Run end-to-end (parse + type-check + interpret + render):
python main.py --run src/tests/fixtures/valid/sample1.rcx
python main.py --run src/tests/fixtures/valid/sample2.rcx
\end{lstlisting}

\noindent At the time of submission the test suite reports \textbf{152/152
green} (\texttt{python -m unittest discover src/tests}). Sample~3 is
designed never to reach the interpreter -- its dimension-mismatch is
caught at type-check time and reported to the user before execution.


\section{Sample 1 -- Pancakes (function + recipe + scale + evaluate)}

\subsection*{Source program}

\begin{lstlisting}
function half(n: int) -> int {
    return n / 2
}

recipe pancakes(servings: int) serves servings {
    ingredient flour : 50 g * servings
    ingredient milk  : 60 ml * servings
    ingredient eggs  : half(servings) * 1 count
    ingredient salt  : 1 pinch

    step "Mix dry" {
        combine(flour, salt)
    }
    step "Add wet" {
        combine(milk, eggs)
    }
    step "Cook batches" at 180 °C for 3 min {
        repeat servings times {
            pour(milk)
            flip()
        }
    }
}

evaluate scale(pancakes(servings: 4), by: 1.5)
\end{lstlisting}

\subsection*{Actual output (rendered by the interpreter)}

\begin{lstlisting}[style=output]
pancakes — serves 6

Ingredients:
  flour: 300 g
  milk: 360 ml
  eggs: 3
  salt: a pinch

Steps:
  1. Mix dry
     - combine(flour, salt)
  2. Add wet
     - combine(milk, eggs)
  3. Cook batches at 180 °C for 3 min
     - pour(milk)
     - flip()
     - pour(milk)
     - flip()
     - pour(milk)
     - flip()
     - pour(milk)
     - flip()
     - pour(milk)
     - flip()
     - pour(milk)
     - flip()
\end{lstlisting}

\subsection*{Discussion}

This program exercises a scalar helper function (\texttt{half}), a
parameterized recipe with a \texttt{serves} expression that depends on
a parameter, four ingredient declarations that mix \texttt{Mass},
\texttt{Volume}, \texttt{Count}, and \texttt{Pinch} types, three step
declarations including one with both \texttt{at} and \texttt{for}
modifiers (in the locked decision~\#23 order), a count-bounded
\texttt{repeat} loop whose count depends on the recipe parameter, and a
\texttt{scale} call at the top level wrapped in an \texttt{evaluate}.
The output shows \texttt{scale(\_, by:~1.5)} correctly multiplying
\texttt{Mass} ($50~\text{g} \times 4 \times 1.5 = 300~\text{g}$) and
\texttt{Volume} ($60~\text{ml} \times 4 \times 1.5 = 360~\text{ml}$)
ingredients and the \texttt{Count} of eggs ($2 \times 1.5 = 3$ after
\texttt{int(round(...))} per plan \S2.6's banker's rounding rule),
while leaving \texttt{Pinch} untouched and the \texttt{at 180~°C}
temperature unchanged -- decision~\#25's ``scale touches quantities
and servings, not Temperature or Duration.''

The six \texttt{pour(milk) / flip()} cycles in step~3 are worth
specific comment because they demonstrate the \textbf{step-body
re-execution rule} locked in D1~\S4. The recipe was instantiated with
\texttt{servings = 4}, then \texttt{scale(\_, by:~1.5)} rebound
\texttt{servings} to~6 in a fresh evaluation environment and
re-executed the step bodies against it. The \texttt{repeat servings
times} loop counts six iterations under the rebound \texttt{servings},
not four. This is the user-intuitive reading of \texttt{scale} --
doubling a recipe doubles the work, not just the header numbers -- and
it is consistent with spec~\S9's ``multiplies the servings field by
the scalar'' once we accept that any step-body construct referring to
\texttt{servings} should reflect the scaled value. Referential
transparency (decision~\#25) is preserved because \texttt{scale} still
produces a fresh \texttt{RtRecipe} value; the re-execution happens
entirely inside the new value's construction, not by mutating the
original.


\section{Sample 2 -- Smoothie (substitution at the call site)}

\subsection*{Source program}

\begin{lstlisting}
recipe smoothie(servings: int) serves servings {
    ingredient milk        : 150 ml * servings
    ingredient banana      : 1 count * servings
    ingredient sweetener   : 10 g * servings
    ingredient oat_milk    : 150 ml * servings

    step "Blend everything" for 1 min {
        combine(milk, banana, sweetener)
        blend()
    }
}

// Vegan variant: substitute happens at the call site, not inside the recipe.
evaluate substitute(smoothie(servings: 2), milk, with: oat_milk, ratio: 1.0)

// Non-vegan variant:
evaluate smoothie(servings: 2)
\end{lstlisting}

\subsection*{Actual output (rendered by the interpreter -- two consecutive evaluates)}

\begin{lstlisting}[style=output]
smoothie — serves 2

Ingredients:
  oat_milk: 300 ml
  banana: 2
  sweetener: 20 g

Steps:
  1. Blend everything for 1 min
     - combine(oat_milk, banana, sweetener)
     - blend()

smoothie — serves 2

Ingredients:
  milk: 300 ml
  banana: 2
  sweetener: 20 g
  oat_milk: 300 ml

Steps:
  1. Blend everything for 1 min
     - combine(milk, banana, sweetener)
     - blend()
\end{lstlisting}

\subsection*{Discussion}

This program demonstrates \texttt{substitute} as a call-site-only
domain operator (decision~\#25, no \texttt{this}) and the bare-IDENT
slot restriction (decision~\#35) that makes the binding-resolution
rule visible at the grammar level. The first \texttt{evaluate} is the
vegan variant: \texttt{substitute(smoothie(servings:~2), milk, with:
oat\_milk, ratio:~1.0)} looks up the binding \texttt{milk} in the
smoothie's ingredient table, replaces it with the pre-declared
\texttt{oat\_milk} ingredient at the given ratio, and produces a
fresh \texttt{RtRecipe} value with \texttt{milk} gone from the
ingredients dict and \texttt{oat\_milk} taking its quantity. The
step-body action \texttt{combine(milk, banana, sweetener)} is
correspondingly rewritten to \texttt{combine(oat\_milk, banana,
sweetener)} because the substitution propagates through every
reference to the substituted ingredient binding -- this is
decision~\#28's ``ingredient identity is the binding, not the
\texttt{name} field'' working as designed.

The second \texttt{evaluate} is the non-vegan variant: \texttt{evaluate
smoothie(servings:~2)} produces a fresh instantiation with no
substitution, so both \texttt{milk} and \texttt{oat\_milk} appear in
the output. The latter is pre-declared in the recipe body precisely so
that the call-site \texttt{substitute} can reference it as a binding;
when no substitute is applied, the pre-declared alternative simply
sits unused alongside the canonical ingredient. This is the
no-\texttt{this} design (spec~\S6, decision~\#25): a recipe cannot
reference itself during construction, so alternative ingredients must
be in scope at the call site, which means they have to be declared
inside the recipe body, which means they appear in the un-substituted
output. The trade-off is the small visual oddity of seeing both
\texttt{milk} and \texttt{oat\_milk} in the non-substituted variant,
accepted in exchange for fully functional, referentially transparent
substitute semantics.


\section{Sample 3 -- Dimension-mismatch type error}

\subsection*{Source program}

\begin{lstlisting}
recipe broken() serves 1 {
    ingredient flour : 200 g
    ingredient water : 100 ml

    step "Combine wet and dry" {
        // ERROR (line below): cannot add Mass and Volume.
        // The implicit ingredient-to-quantity projection turns
        // `flour` into a Mass and `water` into a Volume, then
        // arithmetic fails the dimension-match check.
        let total : Mass = flour + water
    }
}

evaluate broken()
\end{lstlisting}

\subsection*{Type-error trace}

Command run:
\begin{lstlisting}[style=output]
python main.py --typecheck src/tests/fixtures/valid/sample3.rcx
\end{lstlisting}

Verbatim output (exit code 1):

\begin{lstlisting}[style=errmsg]
type error at line 13, col 0: dimension mismatch: cannot + Mass and Volume
\end{lstlisting}

\subsection*{Discussion}

Line~13 of \texttt{sample3.rcx} is the \texttt{let total : Mass = flour
+ water} line. The implicit Ingredient~$\to$~Quantity projection
(decision~\#31) turns \texttt{flour} (\texttt{Ingredient<Mass>}) into a
\texttt{Mass} value and \texttt{water} (\texttt{Ingredient<Volume>})
into a \texttt{Volume} value before the \texttt{+} operator
dispatches. The dimension-match check in \texttt{\_check\_BinaryOp}
then refuses to add two different dimensions and raises
\texttt{TypeCheckError(error\_code=1)} -- spec~\S10 error~\#1
(``dimension mismatch in arithmetic'').

This is the headline-claim demonstration for Recipix: every numeric
quantity carries a dimension at compile time, and any operation that
mixes dimensions is rejected \emph{before execution}. The
dimension-mismatch error is caught entirely by the type checker --
the interpreter never sees \texttt{sample3.rcx} because
\texttt{check()} raises before \texttt{run()} is called. This is what
spec~\S1 means by ``dimensional type discipline'' as the language's
reason to exist: a generic scripting language treats \texttt{200} and
\texttt{100} as raw integers and lets the programmer add grams to
milliliters; Recipix lifts the dimensional structure into the type
system and makes the nonsense impossible to express.

The error message format follows the established \texttt{parse error
at line X, col Y: ...} convention from Part~1 (now \texttt{type error}
at the middle phase), so a user reading the error sees the same shape
regardless of which compiler phase rejected their program. The trace
is pinned by the type checker's \texttt{\_check\_BinaryOp $\to$
\_check\_quantity\_binop} path in \texttt{src/recipix/typechecker.py};
the same path catches any \texttt{Mass + Volume}, \texttt{Temperature
+ Duration}, \texttt{Volume - Mass}, or similar cross-dimension
arithmetic at compile time.


\section{Parse-error tests (carryover from Part 1)}

Five malformed programs from D3~[P1] are retained in
\texttt{src/tests/fixtures/invalid/}. Each demonstrates a parse-error
message with a line number, satisfying the D3~[P1] requirement; they
continue to pass in Part~2 since \texttt{parser.py} was not modified.

\begin{center}
\begin{tabular}{@{}p{0.30\textwidth}p{0.20\textwidth}p{0.45\textwidth}@{}}
\toprule
\textbf{Program} & \textbf{Decision} & \textbf{Parser error} \\
\midrule
\texttt{missing\_brace\_if.rcx}      & \#20 mandatory braces       & \texttt{expected '\{' after 'if' condition (braces are mandatory)} \\
\texttt{wrong\_modifier\_order.rcx}  & \#23 at-before-for          & \texttt{'at' modifier must come before 'for' modifier} \\
\texttt{substitute\_expr\_slot.rcx}  & \#35 bare-IDENT slots       & \texttt{substitute ingredient slot must be a bare identifier} \\
\texttt{nonassoc\_comparison.rcx}    & \#17 non-associative cmp    & \texttt{chained comparison; operators are non-associative} \\
\texttt{quantity\_of\_no\_paren.rcx} & \#34 parens required        & \texttt{quantity\_of requires parenthesized operand} \\
\bottomrule
\end{tabular}
\end{center}

\noindent All five remain green in the parser regression suite under
\texttt{src/tests/test\_parser.py}.


\section{Summary}

Three sample programs cover every control structure (\texttt{if},
\texttt{repeat}, \texttt{foreach} -- the last exercised in the
regression suite), the structured type (\texttt{List<T>} inferred at
literal construction), function and recipe definition and call, all
three domain-specific operators (\texttt{evaluate}, \texttt{scale},
\texttt{substitute}), and the dimensional-correctness type error
caught at compile time. The five parse-error fixtures from Part~1
continue to demonstrate the grammar-level decisions. Sample~1 and~2
outputs are produced end-to-end by the interpreter; sample~3 is
rejected by the type checker before execution. Test count at
submission: \textbf{152/152 green}, zero known failures.

\vfill
\noindent\rule{\textwidth}{0.4pt}
\begin{center}
\small
\emph{End of D3 [P2]. Companion documents: D1 (design specification),
the locked language specification at \texttt{recipix\_v4\_1\_spec.md},
and the joint contract at \texttt{part2\_plan.md}.}
\end{center}

\end{document}
